\documentclass[12pt,a4paper]{article}
\usepackage[utf8]{inputenc}
\usepackage[portuguese]{babel}
\usepackage{geometry}
\usepackage{setspace}
\usepackage{enumitem}
\usepackage{amsmath}
\usepackage{graphicx}

\geometry{margin=2.5cm}
\onehalfspacing

\begin{document}

\begin{center}
\Large \textbf{Questionário Avaliativo}\\
\vspace{0.2cm}
\large \textbf{Disciplina: Revisão Sistemática e Meta-análises}\\
\vspace{0.2cm}
\normalsize Tema 1 – Metodologia de Investigação em Revisão Sistemática
\end{center}

\vspace{0.5cm}

\section*{Instruções}
Este questionário destina-se a estudantes de \textbf{Mestrado e Doutoramento}.  
Assinale \textbf{apenas uma opção correta} em cada questão.

\vspace{0.5cm}

\section*{Questionário}

\begin{enumerate}[label=\textbf{Q\arabic*.}]

\item Qual é o principal objetivo de uma revisão sistemática?
\begin{enumerate}[label=\alph*)]
\item Resumir narrativamente estudos relevantes  
\item Selecionar artigos de acordo com a opinião do investigador  
\item Identificar, avaliar criticamente e sintetizar evidência científica de forma reprodutível  
\item Apresentar resultados preliminares de um estudo experimental  
\end{enumerate}

\item Em investigação científica, a principal diferença entre investigação quantitativa e qualitativa é:
\begin{enumerate}[label=\alph*)]
\item O tipo de revista onde é publicada  
\item O tamanho da amostra  
\item A natureza dos dados e do raciocínio analítico  
\item O uso obrigatório de estatística  
\end{enumerate}

\item Segundo a pirâmide dos níveis de evidência, qual tipo de estudo apresenta maior nível de evidência?
\begin{enumerate}[label=\alph*)]
\item Estudos de caso  
\item Ensaios clínicos não randomizados  
\item Revisões sistemáticas e meta-análises  
\item Opinião de especialistas  
\end{enumerate}

\item Qual das opções abaixo representa um desenho de investigação \textbf{observacional}?
\begin{enumerate}[label=\alph*)]
\item Ensaio clínico randomizado  
\item Estudo de coorte  
\item Estudo experimental in vitro  
\item Simulação computacional  
\end{enumerate}

\item Um elemento fundamental de qualquer estudo científico é:
\begin{enumerate}[label=\alph*)]
\item A confirmação das hipóteses iniciais  
\item A replicabilidade metodológica  
\item O uso obrigatório de software estatístico  
\item A publicação em revista de alto fator de impacto  
\end{enumerate}

\item As chamadas \textbf{medidas de força de associação} incluem:
\begin{enumerate}[label=\alph*)]
\item Média e mediana  
\item Desvio padrão e variância  
\item Odds Ratio, Risk Ratio e Hazard Ratio  
\item Intervalo interquartil  
\end{enumerate}

\item A investigação qualitativa caracteriza-se principalmente por:
\begin{enumerate}[label=\alph*)]
\item Uso exclusivo de dados numéricos  
\item Análise estatística inferencial  
\item Exploração de significados, perceções e contextos  
\item Aplicação de modelos probabilísticos  
\end{enumerate}

\item Qual das seguintes afirmações sobre revisões sistemáticas é correta?
\begin{enumerate}[label=\alph*)]
\item São menos rigorosas que revisões narrativas  
\item Não requerem critérios de elegibilidade  
\item Devem seguir protocolos explícitos e transparentes  
\item Não podem incluir meta-análises  
\end{enumerate}

\item Na avaliação crítica de um artigo científico, uma questão essencial é:
\begin{enumerate}[label=\alph*)]
\item O número de autores do artigo  
\item O país da instituição do autor correspondente  
\item A validade interna e o risco de viés  
\item O fator de impacto da revista  
\end{enumerate}

\item As chamadas \textbf{“10 questões”} na avaliação crítica de artigos científicos visam principalmente:
\begin{enumerate}[label=\alph*)]
\item Avaliar o estilo de escrita científica  
\item Identificar erros tipográficos  
\item Avaliar a robustez metodológica e a credibilidade dos resultados  
\item Determinar a relevância política do estudo  
\end{enumerate}

\end{enumerate}

\vspace{0.7cm}

\section*{Observação Final}
Este questionário pode ser facilmente convertido para \textbf{Google Forms}, utilizando perguntas de múltipla escolha com uma única resposta correta.

\vspace{0.5cm}

\begin{center}
\rule{0.5\linewidth}{0.4pt}\\
\textit{Departamento de Física / Ciências Aplicadas}\\
\textit{Pós-graduação}
\end{center}

\end{document}
